% !TeX root = ../papers/nc-p3.tex
\section{特異な核曲線への母線族の降下と三点繊維の変形}
\label{sec:cofactor-singular-kernel-descent}
\begin{theorem}\label{thm:cofactor-singular-kernel-descent}
定理\ref{thm:cofactor-sectional-kernel}で核像 $K$ が曲線の場合、その次数を $e=24-d$、
正規化 $B$ の種数を $g$ とする。このとき $p_a(K)\ge e-3$ である。
等号なら入力母線族は $K$ 自身からのGrassmann射となり、その射影直線束 $S_K$ からの有限射
$f:S_K\to D$ と射影 $\pi:S_K\to K$ について
\[
                  \operatorname{coker}N=f_*\pi^*\mathcal O_K(1).
\]
特に二十個の印はこの束の長さ三繊維である。$e>3$ なら入力の全母線に会う固定直線はない。
\end{theorem}
\begin{proof}
定理\ref{thm:kernel-curve-scroll}の正規化 $S_B$ と $M\hookrightarrow f_{B*}\pi_B^*\mathcal O_B(1)$
を使う。ここで $\mathcal O_B(1)$ は核曲線の超平面束の引戻しである。
一般入力平面 $H$ で切ると $S_B$ 上では各母線の切断で $B$ 自身となり、
入力平面八次 $E$ の正規化写像 $h:B\to E$ を得る。前記包含の商は $H$ 上で長さ
$a=e-g-3$ を持ち、一般切断ゆえ包含も保たれる。
階数二の $d$ 点で $h$ の繊維長は二である。平面曲線の重複度二は各正規化枝での極大イデアルの位数和だからである。
この二点スキームは $B$ の部分スキームであり、一般 $H$ では $B\to K$ でその長さを保つ。
実際 $K$ の特異点上の有限個の母線と入力の整曲線 $C$ との交点は有限で、
$C$ は直線でない。また異なる二つのパラメータの一般核点が同じであればその核点は特異である。
従ってこの長さ二の上で四つの核座標の評価は全射し、商はそこに寄与しない。
階数三では $M|_E$ は線束であり、核点の写像 $E\dashrightarrow K$ は正則である。
局所的に一つの核座標を単元とすれば、そこでの商長は $E$ の正規化余長に等しい。
滑らかな $K$ の点の付近でこの余長は零である。正規化から $K$ への合成が滑らかな一点での
局所同型であるため、$E$ の局所環もその正規化環と等しくなる。

各特異核点 $z$ 上のパラメータの $E$ における像 $x$ は全て階数三である（上と同じ一般平面の選択）。
完備局所環において
\[
 \widehat{\mathcal O}_{K,z}\ \subseteq     \prod_x\widehat{\mathcal O}_{E,x}\ \subseteq        \prod_{b\mapsto z}\widehat{\mathcal O}_{B,b}.
\]
実際中間環は有限な枝の分割で定まり、左の射は各 $x$ での核点写像で、合成は正規化写像である。
従って $a\le p_a(K)-g$。等号なら各 $z$ で左の包含は同型で、局所環は非自明な冪等元を持たないので
$x$ はただ一つとなる。さらに $B\to E\subset H$ は $K$ 上に降りる。これは各特異点で
射影座標の単元比が完備局所環で降り、正規化の有限余長商への像の消滅は完備化で判定できることによる。
二つの一般平面に対してこれを使う。両平面の共通直線は特異核点の上の有限個の入力母線を避けられるので、
二切断の相異なる入力点の張る直線を取れば $K$ 上の母線射になる。

$S_K$ はこの部分二束の直線束で、$f$ は有限（正規化からの射が有限）。
核座標から $M\hookrightarrow f_*\pi^*\mathcal O_K(1)$ を得る。
入力超平面の相対切断二束は次数八であり、$t$ 次対称冪の階数は $t+1$、次数は $4t(t+1)$。
ゆえに右辺のEuler数は $(t+1)(4t+e+1-p_a(K))=4(t+1)^2$ で等号になる。
生成数と繊維長の一致も従う。
固定直線 $J$ が全母線に会うなら、その交点は動く（固定点なら $N$ はそこで非退化な核曲線全体を消してしまう）。
従って $J$ の一般点は $C$ の外で階数三であり、その核写像が $K$ を支配する。
$J$ 上の六次余因子を原始な核列の平方と共通因子に分けると核写像の次数は三以下であり、$e>3$ では不可能である。
\end{proof}

\begin{proposition}\label{prop:cofactor-kernel-smoothing}
$e>3$ とする。母線写像が $B$ と $K$ の間の整曲線 $K'$ 上に降り、
$p_a(K')=e-3\le3$ なら矛盾する。
\end{proposition}
\begin{proof}[入力二束と基底曲線の滑らかさへの変形]
$K'$ 上で入力四空間 $V_0$ を母線二束 $U$ で割った二束を $F$ とする。
前定理のEuler数の証明を $K'$ と $\mathcal O_K(1)$ の引戻しに用いれば、
ここでも $M=f_*\pi^*\mathcal O_K(1)$（右辺では $K'$ へ引き戻す）である。
従って入力点に対応する $F$ の切断の零スキームの長さはちょうど行列の余階数である。
任意の $z\in K'$ でその入力母線の一般点は整な非直線 $C$ の外なので階数三である。
従ってこの切断の二成分は $\mathcal O_{K',z}$ で極大イデアルを生成する。ゆえに
$K'$ は全点で埋込み次元二以下の局所平面曲線で、双対化層 $\omega$ は可逆である。

Grassmann射の障害群 $H^1(U^*\otimes F)$ は零である。
実際 $\operatorname{Hom}(F,\omega)$ を $V_0\to H^0(\omega)$ の線形写像の空間に入れる。
非零な写像の階数が二以下なら、固定二空間の $F$ への像は一般点で高々一となり
その射影直線が全入力母線に会ってしまう。従って種数二以下ではこのHomは零である。
種数三ではその全ての非零元は三行四列で階数三だから、Hom空間の次元は二以下である
（階数二以下の射影多様体は余次元二で、どの射影平面にも会う）。
よって $F\to U\otimes\omega\subset V_0\otimes\omega$ の四成分が張る空間は高々二次元である。
非零ならその像は固定二空間 $W\subset V_0$ の $\omega$ 倍に入り、
$U$ は一般点で $W$ に会う。これも固定横断線の排除に反する。Serre双対性で消滅が従う。

局所平面曲線を射影的な一次元族で平滑化し、その上でGrassmann射を延長できることを詳しく述べる。
十分高次の非常に豊富な束で $K'$ を射影空間に埋め込めば、局所完全交叉であり、
Euler列から $H^1(T_{\mathbb P}|_{K'})=0$。接・法完全列を特異点での $\mathcal T^1$ まで取ると
$H^1(N_{K'/\mathbb P})=0$ かつ $H^0(N_{K'/\mathbb P})\to H^0(\mathcal T^1)$ は全射である。
ここで $\mathcal T^1$ は有限支持であり、曲線上の第二コホモロジーは消える。
局所完全交叉の方程式の持ち上げの差を法束のコサイクルで補正すれば、
埋込変形から各特異点の変形への写像は形式的に滑らかとなる（局所変形を固定した障害は
$\ker(N\to\mathcal T^1)$ の $H^1$ で、上の接列からこれも零）。
局所平面方程式は一般定数項で平滑化でき、従って有限型の埋込Hilbertスキームで
滑らかな繊維へ向かう曲線を取れる。Grassmann射もアフィン被覆での持ち上げの差を
$H^1(U^*\otimes F)=0$ により補正できるので、相対Homスキームの滑らかさにより
エタール局所でこの族に延長される。$\deg F=8$ と算術種数は一定である。
\end{proof}

\begin{proof}[三点繊維の数]
長さ三部分スキームを印した相対Hilbertスキームと入力 $\mathbb P^3$ の積の上で、
$F^{[3]}(1)$ の普遍切断の零点を取る。中心繊維の零点集合は有限で二十個以上である。
実際有限押出しの生成数が繊維長であり、階数は全点で一以上、階数一点は有限で、
長さ三繊維はそれ自体しか含まない。各中心零点は一般繊維へ変形する。
埋込空間の次元を $r$ とする。長さ三のスキームの各点で埋込み次元は二以下であり、
$\operatorname{Hilb}^3(\mathbb P^r)$ はその点で滑らかな次元 $3r$ の空間である。
実際一般の平面への射影はその長さ三上で閉埋込みとなり、近傍は滑らかな平面の点のHilbertスキーム上の
残り $r-2$ 座標（各々長さ三の関数）のグラフで書ける。
埋込曲線の相対方程式は局所完全交叉により $r-1$ 本であり、長さ三上の消滅で高々 $3(r-1)$ 式。
複数の台点では各台の長さに分けて同じ数となる。さらに二束切断の消滅が六式で、
一次元の平滑化基底を含め零点は局所次元 $3r+3+1-3(r-1)-6=1$ 以上である。
中心繊維が有限なので各点を通る水平成分があり、固有性と上半連続性で一般繊維も有限となる。
従ってその長さは少なくとも中心繊維の相異なる点数である（hensel化で有限射を各中心点に分ければよい）。
滑らかな一般繊維では系\ref{cor:kernel-curve-triple-ruling-genus}と同じ三点対角積分で
長さは $20-4p_a(K')<20$。これは矛盾である。
\end{proof}

\begin{theorem}\label{thm:cofactor-curve-kernel-only-bezout}
$d=18,19,20$ では核像は必ず曲面である。従って整で一般被約な階数二曲線を持つ被約有限窓で
核像が曲線になるのは $d=21$ の七余因子Bezout枝に限られる。
\end{theorem}
\begin{proof}
非退化な整空間曲線のCastelnuovo上界は次数 $e=4,5,6$ で順に $1,2,4$ である。
先の下界 $p_a(K)\ge e-3$ と比較すると $p_a(K)=e-3$ または $(e,p_a)=(6,4)$ だけである。
等号なら第一定理の降下と前命題で矛盾する。

六次種数四では二次式の制限は九次元以下なので二次曲面に載る
（次数十二は標準次数六を超える）。滑らかな二次曲面では種数公式から $(3,3)$、
二次錐でも重み $1,1,2$ の座標における次数六の方程式は通常の三次超平面切断である。
よって $K$ は二次・三次の完全交叉で全点でGorensteinである。
一般入力平面の切断で第一定理の完備環の列に戻ると、今度は左の包含の余長が全点合計で一である。
その非零点は一つの特異点 $z$ である。一次元Gorenstein局所環 $R$ と極大イデアル
$\mathfrak m$ に対して、余長一の有限中間拡大環は存在すればただ一つである。
実際全て $\mathfrak m^{-1}=\operatorname{Hom}_R(\mathfrak m,R)$ に入り、
$0\to\mathfrak m\to R\to k\to0$ の双対で $\mathfrak m^{-1}/R=k$。
この議論は完備化後も成立する。従って平面切断への写像が降りる中間曲線は有限候補しかなく、
いずれも $K$ の構造層を一つの特異点で余長一だけ拡大した部分正規化である。
完備環から元の曲線への移行は、正規化から導手を除いた有限長商の部分代数を取ればよい。

一候補 $K'$ に降りる条件は一般平面のパラメータ上で構成可能である。
具体的には特異点上の有限個のパラメータで消えない入力座標を選び、
母線と平面の交点のアフィン座標をその座標で割って読む。
これらが上記有限長商の指定部分代数に入ることは分母が単元の範囲で代数条件である。
従って一つの $K'$ が稠密開集合上の全平面に使える。二平面との交点を結んで
母線写像自体が $K'$ に降り、$p_a(K')=4-1=3=e-3$。前命題を適用すると矛盾する。
\end{proof}
